\section{Introduction}
\label{sec:introduction}

The promise of algorithmic redistricting rests on a claim of structural
impartiality: if the algorithm cannot see partisan data, the maps it draws
cannot be deliberately gerrymandered.
This claim is widely accepted in the literature~\citep{chen2013unintentional,
browdy1990computer, altman2011bard} and is embedded in the logic of proposed
legislative remedies such as the \dia{}.

A sophisticated objection, however, runs as follows: even if the \emph{algorithm}
cannot use partisan data, the \emph{actors who deploy the algorithm} can still
game the outcome by manipulating the inputs, parameters, geographic definitions,
or presentation of results.
Algorithmic neutrality of the core computation is not sufficient if the surrounding
context allows a hostile actor to steer outcomes.

This paper takes the objection seriously.
Rather than dismissing it, we systematically enumerate every gaming vector
available to a hostile actor who:
\begin{enumerate}
\item Knows the \bisect{} algorithm in full (the algorithm is open-source);
\item Controls the parameters passed to the algorithm;
\item Has access to the input data pipeline (Census files, shapefiles);
\item Decides which results to present to courts, legislatures, and the public;
\item Can choose which version of the algorithm to run.
\end{enumerate}

For each gaming vector, we answer three questions: (a)~what is the attack?
(b)~how much partisan shift can the attack produce? and (c)~what is the
detection mechanism that catches the attack?

\subsection{Main Findings}

The central finding is that all plan-affecting gaming vectors combined produce
at most 0.5 Democratic seats of variation nationally under additive worst-case
assumptions, and at most 0.3 seats in practice (Section~\ref{sec:combined}).
This is comparable to the noise floor established by B.7's seed-variance analysis
(coefficient of variation $<$ 2\%) and is far smaller than the partisan
effects of algorithm \emph{structure} documented in B.0 (approximately
18 seats in a worst-case comparison of algorithm families).

The finding is significant for two reasons.
First, it shows that parameter tuning within a fixed algorithm is not a
meaningful gaming vector---a hostile actor who knows the bisect algorithm
cannot use parameter choices to manufacture a meaningfully partisan outcome
while maintaining the appearance of algorithmic neutrality.
Second, it clarifies where the real design choices lie: in algorithm structure
(which B.0 shows should be pre-specified in statute) and geographic definitions
(which F.3's threshold rule constrains).

\subsection{The DIA Response}

The \dia{} protocol converts each gaming vector from a threat to a
formality~(Section~\ref{sec:dia}).
Parameters are fixed in a SHA-256 formula before Census data are released,
eliminating post-data optimization.
Input data integrity is verified by hashing against Census Bureau-published
checksums.
Geographic definitions are specified by statute and locked into the manifest.
Result presentation follows a mandatory disclosure standard.
Algorithm version is pinned in the reproducibility manifest and must match
the SHA-256 of the deployed binary.

Under the \dia{} protocol, no gaming vector can produce an undetectable
partisan outcome.
Detected deviations from the manifest are automatic grounds for plan rejection.

\subsection{Track R Structure}

This paper (R.0) provides the taxonomy and establishes the maximum effect
bounds.
Subsequent papers treat each major plan-affecting vector in detail:
R.1 analyzes parameter gaming using B.17's 50-state sweep;
R.2 analyzes input data manipulation and the Census P.L.\ 94-171 integrity chain;
R.3 analyzes geographic definition gaming and resolution choices;
R.4 analyzes the audit chain as a formal gaming defense, including
cryptographic analysis and Daubert-standard reproducibility.

\subsection{Road Map}

Section~\ref{sec:taxonomy} introduces the five-vector taxonomy.
Sections~\ref{sec:v1}--\ref{sec:v5} treat each vector with attack
description, effect measurement, and defense.
Section~\ref{sec:combined} quantifies the joint maximum effect.
Section~\ref{sec:dia} describes the \dia{} countermeasures.
Section~\ref{sec:conclusion} concludes.
